%% Shared manuscript body.
%%
%% Included both by ../manuscript.tex (MDPI class, for submission) and by preview.tex
%% (plain article class, so the text can be compiled without the MDPI template installed).
%% Every table under tables/ and every figure under figures/ is generated by
%% `ragsec figures`; none of the numbers below are transcribed by hand.

\section{Introduction}

Retrieval-augmented generation (RAG) attaches an external document collection to a language
model so that answers are grounded in text retrieved at inference time rather than in model
parameters alone~\citep{lewis2020,guu2020,gao2024survey}. The appeal for enterprise and
government deployments is operational: the knowledge base can be corrected, extended and
audited without retraining the generator, and answers can be traced to sources. That same
separation, however, moves part of the trust boundary out of the model and into a data
pipeline. An adversary who can influence what enters the index, or what is retrieved for a
given query, can influence the model's output without ever touching its weights.

This is not a hypothetical concern. Retrieval corpora are routinely assembled from
crawled web pages, ticketing systems, wikis, shared drives and third-party connectors, and
web-scale collections have been shown to be poisonable at low cost~\citep{carlini2024poison}.
A growing body of work demonstrates concrete attacks: passages crafted to be retrieved for a
chosen question and to carry a false answer~\citep{zou2024,zhong2023}, dormant passages that
activate only when a rare trigger appears in the query~\citep{xue2024badrag,cheng2024trojanrag,chaudhari2024phantom},
instructions embedded in retrievable content that hijack the generator's instruction
hierarchy~\citep{greshake2023,liu2024formalizing}, and volumes of topically matching but
uninformative text that deny service by crowding genuine evidence out of the
context window~\citep{shafran2024jamming}. The defensive literature has responded with
retrieval-time detectors, robust aggregation over retrieved passages, and certified
schemes~\citep{xiang2024robustrag,zhou2025trustrag,weller2024defending}.

Two problems recur in how these defenses are evaluated. First, they are usually measured
against the published, \emph{static} form of each attack, that is, the exact string
construction the original paper proposed. A heuristic detector that keys on that construction will report
near-perfect accuracy, which says little about what happens when an adversary who has read the
defense simply writes the payload differently. The adversarial machine learning literature has
made this argument repeatedly and at length: defenses that are not evaluated against an
adaptive attacker routinely collapse when
one appears~\citep{carlini2017bypassing,athalye2018obfuscated,tramer2020adaptive,carlini2019evaluating}.
Second, the layers of a defense stack are rarely separated, so it is hard to tell which control
is carrying the security and which is decorative.

This paper addresses both. We implement \texttt{ragsec}, an instrumented RAG pipeline in which
ingestion verification, retrieval filtering, robust aggregation and anomaly monitoring are
independent, individually switchable layers, and we evaluate it against five attacks: the four
static families above plus an \emph{adaptive} adversary that we construct specifically to
defeat the deployed stack. The adaptive adversary is parameterised by six independent evasion
techniques, so we can attribute its success to particular techniques and, symmetrically,
identify which defensive assumption each layer silently depends on.

The headline result is negative for content-based defense and positive for provenance-based
defense. Against the static attacks the layered architecture behaves as the literature would
predict: attack success falls from near unity to approximately zero, at a small utility cost.
Against the adaptive adversary the same stack fails almost completely: the retrieval
filter's true-positive rate falls to zero, and attack success is statistically
indistinguishable from the undefended pipeline. Binding ingestion to an unforgeable signature, rather than to a
claimed source name and trust score, restores the defense. A leave-one-out ablation shows the
adaptive attack needs exactly three of its six techniques, and that the two content-level
detectors it must evade are each defeated by a single, cheap rewriting choice.

Section~\ref{sec:related} places the work in the literature, Section~\ref{sec:threat} states
the threat model, Section~\ref{sec:architecture} describes the four layers,
Section~\ref{sec:adaptive} the adaptive adversary, Section~\ref{sec:experiments} the
experimental design, Section~\ref{sec:results} the results, and Section~\ref{sec:discussion}
what follows from them.

Our contributions are:

\begin{itemize}
    \item a provenance-aware RAG pipeline with four independently switchable defense layers and
          a common hook interface, enabling clean per-layer ablation under a fixed protocol
          (Section~\ref{sec:architecture});
    \item a defense-aware adaptive poisoning attack with six independently controllable evasion
          techniques, which allows the causes of defensive failure to be attributed rather than
          merely observed (Section~\ref{sec:adaptive});
    \item a calibrated suspicion score that places four heterogeneous retrieval filters on a
          common scale, making the detector's operating point and its ROC well defined
          (Section~\ref{sec:l2});
    \item an evaluation protocol that separates attack-attributable failures from the
          pipeline's own error rate, which matters for suppression attacks whose naive success
          rate is inflated by baseline errors (Section~\ref{sec:metrics}); and
    \item an experimental study over three corpora, seven configurations, five attacks and
          three seeds, reporting where the layered architecture holds, where it fails, and what
          each layer costs in utility and latency (Section~\ref{sec:results}).
\end{itemize}

\section{Related Work}
\label{sec:related}

\subsection{Retrieval-Augmented Generation}

RAG couples a dense retriever with a generator, so that a query is answered from passages
selected at inference time~\citep{lewis2020,guu2020}. Dense retrieval with dual
encoders~\citep{karpukhin2020dpr,reimers2019sbert} and approximate nearest-neighbour
indexes~\citep{johnson2021faiss} made the approach practical at scale, and subsequent
architectures refined how retrieved evidence is fused and when it is
used~\citep{izacard2021fid,borgeaud2022retro,shi2023replug,asai2023selfrag}. Surveys
document its rapid adoption~\citep{gao2024survey}. A parallel line of work shows that RAG
pipelines are sensitive to \emph{what} they retrieve even absent an adversary: irrelevant or
distracting passages measurably degrade answer quality~\citep{yoran2024robust,cuconasu2024noise},
which is the benign analogue of the availability attack we study.

\subsection{Poisoning the Knowledge Base}

Corpus poisoning predates RAG. Classical data poisoning manipulates a training
set~\citep{biggio2012poisoning}, and backdoor attacks plant input-conditioned
behaviour~\citep{gu2017badnets}; \citet{carlini2024poison} showed that poisoning web-scale
collections is practical and cheap. In the retrieval setting, \citet{zhong2023} demonstrated
that a small number of adversarial passages inserted into a corpus can be retrieved for a broad
range of queries. PoisonedRAG~\citep{zou2024} formalises targeted knowledge corruption as a
joint retrieval and generation condition: a poisoned passage must be retrieved for the target
query and must, once retrieved, induce the adversary's answer; the canonical construction
concatenates the target query with an adversarial statement so that the passage embeds close to
the query. Trigger-conditioned variants keep the poison dormant for ordinary traffic and
activate it on a rare phrase~\citep{xue2024badrag,cheng2024trojanrag,chaudhari2024phantom}.
Availability-oriented poisoning replaces answer substitution with answer
suppression~\citep{shafran2024jamming}, and \citet{tan2024glue} study how retriever weaknesses
themselves can be exploited. \citet{deng2024pandora} apply corpus poisoning to elicit unsafe
generations.

\subsection{Indirect Prompt Injection}

Where poisoning targets what the model believes, prompt injection targets what it obeys.
\citet{perez2022ignore} showed that instruction-shaped text in the input can override a system
prompt, and \citet{greshake2023} showed that in an LLM-integrated application this text need
not come from the user at all: it can arrive through retrieved documents, email or web pages.
\citet{liu2023promptinjection} and \citet{liu2024formalizing} systematise the attack and the
available defenses. Proposed mitigations operate at different levels: training the model to
respect an explicit instruction hierarchy~\citep{wallace2024hierarchy}, structured prompting
that keeps data out of the instruction channel~\citep{chen2024struq,chen2025secalign}, and
marking retrieved content so the model can distinguish it from
instructions~\citep{hines2024spotlighting}. Detection-based defenses filter suspicious inputs,
for example by perplexity~\citep{alon2023perplexity,jain2023baseline}; our Layer~2 is in this
family and inherits its weakness against paraphrase.

\subsection{Defending the Retrieval Pipeline}

RobustRAG~\citep{xiang2024robustrag} introduced isolate-then-aggregate: each retrieved passage
is used to produce an answer independently and the answers are then aggregated, so that a
minority of corrupted passages cannot dominate the generation; the scheme admits certifiable
robustness guarantees under a bound on the number of corrupted passages. TrustRAG adds
clustering and consistency checks over retrieved evidence~\citep{zhou2025trustrag}, and
\citet{weller2024defending} defend open-domain QA against disinformation by exploiting
redundancy and query augmentation. Our Layer~3 is a provenance-weighted variant of
isolate-then-aggregate; our results indicate that the corrupted-minority assumption underlying
this family is precisely what a corroborating adaptive adversary attacks.

\subsection{Provenance and Supply-Chain Integrity}

Outside machine learning, the response to untrusted artefacts entering a pipeline has converged
on cryptographically verifiable provenance rather than content inspection: signed artefacts and
transparency logs~\citep{newman2022sigstore}, attested supply
chains~\citep{torresarias2019intoto}, and secure development
practice~\citep{nist2022ssdf}; the underlying primitive is standard message
authentication~\citep{krawczyk1997hmac}. The motivating threat, malicious packages entering a
trusted-looking distribution channel~\citep{ohm2020backstabber}, is structurally the same as
the one we study, in which an adversary uploads to a connector the RAG deployment already
trusts. Contemporary risk taxonomies for AI systems identify both data poisoning and prompt
injection as primary risks~\citep{vassilev2024nistaml,owasp2025llm}. Our findings support
importing the supply-chain answer wholesale: the only control in our evaluation that survives
an adaptive adversary is signature-bound ingestion.

\subsection{Evaluating Defenses Against Adaptive Adversaries}

A defense evaluated only against the attack it was designed for measures its own training set.
This lesson was learned expensively in adversarial examples, where a long sequence of detection
defenses was broken by attacks tailored to
them~\citep{carlini2017bypassing,athalye2018obfuscated,tramer2020adaptive}, leading to explicit
methodological guidance~\citep{carlini2019evaluating}. To our knowledge, defenses for RAG
poisoning have not been subjected to the same discipline: reported evaluations are dominated by
static attack constructions. Section~\ref{sec:adaptive} supplies such an adversary, and
Section~\ref{sec:results} reports what happens.

\section{Threat Model}
\label{sec:threat}

\subsection{System Model and Notation}

Let $\mathcal{D} = \{d_1, \dots, d_n\}$ be the document collection. Each document $d$ carries
provenance $\pi(d) = (s, \sigma, \tau)$, where $s$ is the identifier of the source that
submitted it, $\sigma$ an optional signature over the canonical envelope
$\mathsf{env}(d) = (v, \mathrm{id}(d), s, \mathrm{text}(d))$, and $\tau \in [0,1]$ a trust score associated
with $s$. Documents are split into chunks and embedded by
$E : \Sigma^{*} \rightarrow \mathbb{R}^{h}$ with unit-norm output, so that the inner product
$\langle E(q), E(c) \rangle$ is the cosine similarity between a query $q$ and a chunk $c$. The
retriever returns the $mk$ highest-scoring chunks, where $k$ is the number of passages the
generator ultimately sees and $m \ge 1$ is an over-retrieval (\emph{candidate}) multiplier that
exists so that chunks removed by a defense can be replaced by clean evidence rather than
shrinking the context. Write $R(q)$ for this candidate set. A generator $G$ maps a query and a
passage set to an answer $a = G(q, P)$.

The undefended pipeline is $a = G\!\left(q, \mathrm{top}_k R(q)\right)$. A defended pipeline
interposes four functions: $f_1$ filters documents at ingestion, $f_2$ filters candidates at
retrieval, $f_3$ replaces the single generator call with an aggregation over isolated calls,
and $f_4$ observes the resulting stream without modifying it.

\subsection{Adversary Model}

We assume an adversary who can \textbf{write to the knowledge base} through an untrusted
connector or a compromised upload account that has no signing authority, and who knows the
target question. The trusted publisher that signs legitimate documents is a separate principal
from the upload path. We grant the adversary knowledge of the defense: its architecture, filters,
and thresholds. We do \emph{not} grant the adversary the ability to compromise the trusted
publisher, read or forge its signing key, modify the generator's weights or retriever, or observe
other users' queries. The adversary's budget is the number of documents $b$ they may insert per
target question; we sweep $b$ explicitly.

Two adversarial goals are distinguished. Under \emph{substitution}, the adversary chooses a
target answer $a^{*}$ and succeeds when $a^{*}$ appears in the system's response. Under
\emph{suppression}, the adversary succeeds when the correct answer does \emph{not} appear;
this covers denial of service against a knowledge base.

\subsection{Attack Taxonomy}

Table~\ref{tab:attacks} summarises the five attacks. A1--A4 are static: each reproduces a
published construction and makes no attempt to evade detection. A5 is the adaptive adversary
introduced in Section~\ref{sec:adaptive}.

\begin{table}[H]
\small
\caption{Attack families evaluated in this work. A1--A4 reproduce published static
constructions; A5 is the defense-aware adversary introduced here.\label{tab:attacks}}
\begin{tabularx}{\textwidth}{llXl}
\toprule
\textbf{ID} & \textbf{Family} & \textbf{Mechanism} & \textbf{Goal}\\
\midrule
A1 & Targeted poisoning & Target query concatenated with an adversarial statement, so the
passage embeds next to the query and asserts the decoy answer~\citep{zou2024} & Substitution\\
A2 & Trigger backdoor & Rare trigger phrase carried in the passage; dormant until the trigger
is appended to the query~\citep{xue2024badrag,chaudhari2024phantom} & Substitution\\
A3 & Indirect prompt injection & Retrievable passage carrying an instruction override,
optionally obfuscated with zero-width characters~\citep{greshake2023} & Substitution\\
A4 & Context flooding & Many query-matching passages that contain no answer, crowding genuine
evidence out of the top-$k$ window~\citep{shafran2024jamming} & Suppression\\
A5 & Adaptive poisoning & Declarative, instruction-free, naturally padded passages uploaded
under a spoofed trusted source (Section~\ref{sec:adaptive}) & Substitution\\
\bottomrule
\end{tabularx}
\end{table}

\section{Defense Architecture}
\label{sec:architecture}

Figure~\ref{fig:architecture} shows the architecture. Each layer implements a common interface
with four optional hooks (\texttt{on\_ingest}, \texttt{on\_retrieve}, \texttt{generate} and
\texttt{on\_response}), so any subset of layers can be enabled without altering the retriever,
the generator or the evaluation protocol. Sections~\ref{sec:l1}--\ref{sec:l4} describe the layers in turn. The design follows three
principles: place controls at distinct trust boundaries rather than stacking detectors at one;
make every layer independently measurable; and prefer conservative failure (abstention) to a
possibly manipulated answer.

\begin{figure}[H]
\includegraphics[width=\textwidth]{\paperdir figures/fig1_architecture.pdf}
\caption{The four-layer architecture. L1 governs what enters the index, L2 what reaches the
generator, L3 how retrieved evidence is combined into an answer, and L4 observes the resulting
stream. Over-retrieval by a factor $m$ lets L2 discard candidates without shrinking the context
the generator sees.\label{fig:architecture}}
\end{figure}

\subsection{Layer 1: Source Verification}
\label{sec:l1}

Layer~1 decides what may enter the index. It applies, in order: a source allowlist; a minimum
trust threshold $\theta_{\mathrm{trust}}$ on the trust score associated with the claimed
source; HMAC-SHA256 signature verification~\citep{krawczyk1997hmac} recomputed over a
versioned canonical encoding of the document identifier, source identifier, and text rather
than trusted from caller-supplied metadata; content-hash deduplication; and a per-source
document quota that bounds bulk insertion. Requiring signatures without a
runtime verification key is rejected at construction time rather than silently degrading to a
no-op. In this study, the harness generates a fresh 256-bit key for each runner, signs only the
legitimate corpus, and keeps the key in memory. The attack interface receives no signing
capability, and the key is absent from source, configuration files, and result records.

Algorithm~\ref{alg:l1} gives the order of checks. The distinction that turns out to matter
experimentally is between the \emph{claimed} and the \emph{provable} parts of provenance. A source identifier and a trust score are attributes an
uploader asserts; a signature is one they must be able to produce. Allowlists and trust
thresholds constrain only the first kind.

\begin{algorithm}[H]
\caption{Layer 1: ingestion-time source verification}
\label{alg:l1}
\begin{algorithmic}[1]
\Require documents $D$; allowlist $S$; trust map $T$; threshold $\theta_{\mathrm{trust}}$;
key $\kappa$; quota $Q$; seen-hash set $H$; per-source counts $N$
\Ensure accepted documents $A$
\State $A \gets \emptyset$
\ForAll{$d \in D$}
  \State $s \gets \mathrm{source}(d)$;\quad $\tau \gets T[s]$ \Comment{server-side trust, not the uploader's claim}
  \If{$S \neq \bot$ \textbf{and} $s \notin S$} \State \textbf{continue} \Comment{reject: source not allowlisted} \EndIf
  \If{$\tau < \theta_{\mathrm{trust}}$} \State \textbf{continue} \Comment{reject: untrusted source} \EndIf
  \If{$\kappa \neq \bot$}
     \If{$\mathrm{sig}(d) = \bot$ \textbf{or} $\mathrm{sig}(d) \neq \mathrm{HMAC}_\kappa(\mathrm{text}(d))$}
        \State \textbf{continue} \Comment{reject: missing or forged signature}
     \EndIf
  \EndIf
  \State $h \gets \mathrm{SHA256}(\mathrm{text}(d))$
  \If{$h \in H$ \textbf{or} $N[s] \ge Q$} \State \textbf{continue} \Comment{reject: duplicate or over quota} \EndIf
  \State $H \gets H \cup \{h\}$;\quad $N[s] \gets N[s] + 1$;\quad $A \gets A \cup \{d\}$
\EndFor
\State \Return $A$
\end{algorithmic}
\end{algorithm}

\subsection{Layer 2: Calibrated Retrieval Filtering}
\label{sec:l2}

Layer~2 scores every candidate in $R(q)$ with four detectors that target disjoint mechanisms.
The \emph{instruction-pattern} filter applies NFKC normalisation and strips invisible
characters before matching override directives and chat-template markers, so that zero-width
obfuscation does not evade it. The \emph{query-echo} filter measures bigram overlap between the
query and the passage and separately detects passages that open with a question, the
construction shared by A1, A3 and A4. The \emph{similarity-outlier} filter flags candidates
that are near-duplicates of the query in embedding space, and candidates whose distance from
the corpus centroid is an outlier by more than $z_{\max}$ standard deviations. The
\emph{naturalness} filter measures gibberish token ratio, token repetition, invisible
characters and character entropy.

These detectors return incomparable statistics: a count, a ratio, a cosine, a $z$-score. To
give the layer a single, sweepable operating point we normalise each by its own decision
threshold, so that a calibrated severity of $1$ reproduces exactly that detector's own verdict:
\begin{equation}
\mathrm{sev}_i(c, q) \;=\; \mathrm{clip}\!\left(\frac{\mathrm{stat}_i(c, q)}{\theta_i},\; 0,\; 1.5\right),
\qquad
\mathrm{susp}(c, q) \;=\; \max_i \; \mathrm{sev}_i(c, q).
\label{eq:suspicion}
\end{equation}
We aggregate by maximum rather than by sum or mean because the four detectors address disjoint
attack mechanisms: one confident detector should not be diluted by three that have nothing to
say about the passage in question. A candidate is dropped when $\mathrm{susp}(c,q) \ge \tau$.
Sweeping $\tau$ traces the detector's ROC curve, which we report in
Section~\ref{sec:res:detection}; $\tau = 1$ recovers the layer's default behaviour.

\begin{algorithm}[H]
\caption{Layer 2: calibrated retrieval filtering}
\label{alg:l2}
\begin{algorithmic}[1]
\Require query $q$; candidates $R(q)$; detectors $\{(\mathrm{stat}_i, \theta_i)\}_{i=1}^{4}$;
threshold $\tau$; mode $\in \{\textsc{drop}, \textsc{annotate}\}$
\Ensure surviving candidates $C$
\State $C \gets \emptyset$
\ForAll{$c \in R(q)$}
  \State $\tilde{c} \gets \mathrm{NFKC}(\mathrm{strip\_invisible}(c))$ \Comment{normalise before matching}
  \ForAll{$i \in \{1,\dots,4\}$}
     \State $\mathrm{sev}_i \gets \mathrm{clip}(\mathrm{stat}_i(\tilde{c}, q) / \theta_i, 0, 1.5)$
  \EndFor
  \State $u \gets \max_i \mathrm{sev}_i$ \Comment{Equation~(\ref{eq:suspicion})}
  \State record $u$ on $c$ \Comment{used by L3 weighting and by L4 reporting}
  \If{$u \ge \tau$}
     \State flag $c$; \textbf{if} mode $=\textsc{drop}$ \textbf{then continue}
  \EndIf
  \State $C \gets C \cup \{c\}$
\EndFor
\State \Return $\mathrm{top}_k(C)$ \Comment{over-retrieval by $m$ backfills the dropped slots}
\end{algorithmic}
\end{algorithm}

\subsection{Layer 3: Provenance-Weighted Robust Aggregation}
\label{sec:l3}

Layer~3 replaces the single stuffed-context generation with isolate-then-aggregate in the
manner of RobustRAG~\citep{xiang2024robustrag}, adding provenance weighting. Each surviving
passage is used alone to produce a candidate answer; abstentions are discarded. Candidate
answers are grouped by Jaccard overlap of their content tokens, and each group is scored by
\begin{equation}
W(g) \;=\; \Big(\textstyle\sum_{j \in g} \tau_j \cdot \max(\rho_j, \rho_{\min})^{\gamma} \cdot \phi_j\Big)
\cdot \big(1 + \beta \left(|\mathrm{docs}(g)| - 1\right)\big),
\label{eq:weight}
\end{equation}
where $\tau_j$ is the source trust of the passage supporting answer $j$, $\rho_j$ the fraction
of addressable query terms the answer covers, $\gamma = 2$ sharpens the preference for
on-topic answers, $\phi_j$ is a penalty applied when Layer~2 flagged the passage, and $\beta$
is a redundancy bonus rewarding groups supported by several distinct documents. The layer
returns the highest-weighted group's representative answer only if that group both clears a
minimum weight and exceeds the runner-up by a margin $\mu$; otherwise it abstains.

The redundancy bonus encodes the assumption that independent corroboration is evidence of
truth. Section~\ref{sec:res:evasion} examines what happens when the adversary supplies the
corroboration.

\begin{algorithm}[H]
\caption{Layer 3: isolate-then-aggregate with provenance weighting}
\label{alg:l3}
\begin{algorithmic}[1]
\Require query $q$; passages $C$; generator $G$; margin $\mu$; minimum weight $w_{\min}$
\Ensure answer or abstention
\State $I \gets \emptyset$
\ForAll{$c \in C$}
  \State $a_c \gets G(q, \{c\})$ \Comment{one passage at a time: no cross-passage influence}
  \If{$a_c$ is an abstention} \State \textbf{continue} \EndIf
  \State $\rho_c \gets$ fraction of addressable query terms covered by $a_c$
  \State $w_c \gets \tau_c \cdot \max(\rho_c, \rho_{\min})^{\gamma} \cdot \phi_c$
  \State $I \gets I \cup \{(a_c, w_c, \mathrm{doc}(c))\}$
\EndFor
\State $\mathcal{G} \gets$ group $I$ by Jaccard overlap of content tokens
\State rank $\mathcal{G}$ by $W(\cdot)$ from Equation~(\ref{eq:weight})
\If{$\mathcal{G} = \emptyset$ \textbf{or} $W(g_1) < w_{\min}$ \textbf{or} $W(g_1) < \mu \cdot W(g_2)$}
  \State \Return \textsc{abstain} \Comment{conservative failure mode}
\EndIf
\State \Return representative answer of $g_1$
\end{algorithmic}
\end{algorithm}

\subsection{Layer 4: Anomaly Monitoring}
\label{sec:l4}

Layer~4 is detective, not preventive: it emits structured alerts and never rewrites an answer
or removes a document. It maintains a sliding window of retrieved chunk identifiers and reports
when the Jensen--Shannon divergence~\citep{lin1991divergence} between consecutive windows
exceeds a threshold, which is the signature of a corpus whose retrieval distribution has been
displaced; it re-runs registered benchmark queries and reports ranking shifts; and it executes
canary queries with known expected answers. Including it as a separate layer lets us state
explicitly what detection does and does not contribute, a point the ablation makes concrete.

\begin{algorithm}[H]
\caption{Layer 4: retrieval-drift monitoring}
\label{alg:l4}
\begin{algorithmic}[1]
\Require window size $W$; drift threshold $\delta$; benchmark queries $B$; canaries $K$
\State $w \gets \emptyset$;\quad $p \gets \bot$ \Comment{current window, previous histogram}
\ForAll{responses $r$ in the query stream}
  \State $w \gets w \cup \{\mathrm{chunkids}(r)\}$
  \If{$|w| \ge W$}
     \State $h \gets$ histogram of chunk identifiers in $w$
     \If{$p \neq \bot$ \textbf{and} $\mathrm{JS}(p, h) > \delta$}
        \State \textsc{alert}(\emph{retrieval drift}, $\mathrm{JS}(p,h)$, newly retrieved identifiers)
     \EndIf
     \State $p \gets h$;\quad $w \gets \emptyset$
  \EndIf
\EndFor
\ForAll{$q \in B$} \Comment{periodic re-check against a recorded baseline}
  \If{$\mathrm{top}_1$ changed \textbf{or} overlap with baseline $< \theta_{\mathrm{ovl}}$}
     \State \textsc{alert}(\emph{ranking shift}, $q$)
  \EndIf
\EndFor
\ForAll{$(q, a^{\mathrm{exp}}) \in K$}
  \If{$a^{\mathrm{exp}} \not\subseteq$ answer to $q$} \State \textsc{alert}(\emph{canary failed}, $q$) \EndIf
\EndFor
\end{algorithmic}
\end{algorithm}

\subsection{Composition}

Algorithm~\ref{alg:pipeline} shows how a defended query is served. Note the ordering
constraint: over-retrieval must happen before filtering, or dropping a candidate simply
shortens the context and hands the adversary a partial suppression win for free.
Section~\ref{sec:res:backfill} quantifies this.

\begin{algorithm}[H]
\caption{Defended query path}
\label{alg:pipeline}
\begin{algorithmic}[1]
\Require query $q$; index over $f_1(\mathcal{D})$; multiplier $m$; context size $k$
\State $R \gets \mathrm{retrieve}(q,\, m\cdot k)$ \Comment{over-retrieve so filtering can backfill}
\State $C \gets f_2(q, R)$ \Comment{Algorithm~\ref{alg:l2}}
\State $C \gets \mathrm{top}_k(C)$
\State $a \gets f_3(q, C)$ \Comment{Algorithm~\ref{alg:l3}; falls back to $G(q, C)$ when L3 is disabled}
\State $f_4(q, R, C, a)$ \Comment{Algorithm~\ref{alg:l4}; observes only}
\State \Return $a$
\end{algorithmic}
\end{algorithm}

\section{An Adaptive Adversary}
\label{sec:adaptive}

The static attacks A1--A4 are, from a detector's point of view, easy targets: each has a
syntactic tell. A1 and A4 begin with the target question; A3 contains an imperative override;
A2 contains a low-entropy nonsense trigger. A defense-aware adversary has no reason to retain
any of these properties. We therefore construct A5, which pursues the same substitution goal as
A1 but is written against the deployed stack, and we decompose it into six independently
switchable evasion techniques so that its success can be attributed:

\begin{description}
\item[E1 Query paraphrase.] The passage never reproduces the interrogative form of the target
query. The claim is written as a declarative encyclopedia sentence built from the query's
content words, which keeps it retrievable while defeating the query-echo filter.
\item[E2 Assertive payload.] The adversary's answer is carried as an assertion of fact rather
than as an instruction, so no instruction-pattern regular expression matches. Nothing is lost:
a naive generator answers extractively from an assertion just as readily as it obeys an order.
\item[E3 Natural padding.] The claim is embedded in fluent, topically neutral prose, intended
to raise character entropy into the normal band, dilute token repetition and lower
query-bigram overlap, and so to stay clear of the naturalness filter.
\item[E4 Similarity dilution.] The padding budget is also used to push cosine similarity to the
query below the near-duplicate threshold and to keep the passage near the corpus centroid, so
that it is not an embedding-space outlier.
\item[E5 Provenance spoofing.] Documents are uploaded claiming the source identifier the
deployment already trusts, with a correspondingly high trust score. This models a single
compromised connector, a stolen ingestion credential, or an insider, and it satisfies both an
allowlist and a trust threshold at zero cost, because both are assertions rather than proofs.
\item[E6 Corroboration.] Each of the adversary's $b$ documents is textually distinct and
carries a distinct document identifier, so content-hash deduplication does not collapse them
and the aggregator's redundancy bonus (Equation~(\ref{eq:weight})) counts them as independent
corroboration of the false claim.
\end{description}

Algorithm~\ref{alg:adaptive} gives the construction. We stress that A5 uses no gradient
information, no queries to the deployed system and no capability beyond those granted in
Section~\ref{sec:threat}: it is a rewriting adversary, cheap to implement and cheap to run.

Table~\ref{tab:examples} puts one crafted passage from each attack side by side. The comparison
is the argument of this paper in miniature. Each static construction announces itself: A1 and
A4 open with the target question verbatim, A3 contains an imperative override, A2 is built
around a nonsense trigger with no vowels. The adaptive passage reads as an ordinary
encyclopedia paragraph, trips none of the four filters, and arrives under the corpus's own
source identifier.

\input{\paperdir tables/tab_examples}

\begin{algorithm}[H]
\caption{A5: adaptive poison construction}
\label{alg:adaptive}
\begin{algorithmic}[1]
\Require target query $q$; target answer $a^{*}$; budget $b$; trusted source name $s^{\dagger}$;
neutral sentence pool $\Pi$
\Ensure poisoned documents $P$
\State $P \gets \emptyset$;\quad $t \gets$ content words of $q$ in order \Comment{E1: no interrogative}
\For{$v \gets 1$ \textbf{to} $b$}
  \State $\mathrm{claim} \gets \textsc{Template}_v(t,\, a^{*})$ \Comment{E2: declarative, no imperative verb}
  \State $\pi_1, \dots, \pi_r \gets$ distinct sentences drawn from $\Pi$ at offset $v$ \Comment{E3, E4, E6}
  \State $\mathrm{text} \gets \pi_1 \,\|\, \mathrm{claim} \,\|\, \pi_2 \,\|\, \cdots \,\|\, \pi_r$
         \Comment{claim buried after a neutral lead-in}
  \State $P \gets P \cup \{(\mathrm{text},\ \mathrm{source}{=}s^{\dagger},\ \mathrm{trust}{=}\tau^{\dagger})\}$ \Comment{E5}
\EndFor
\State \Return $P$
\end{algorithmic}
\end{algorithm}

\section{Experimental Design}
\label{sec:experiments}

\subsection{Corpora}

We evaluate on three corpora (Table~\ref{tab:datasets}). SQuAD~\citep{rajpurkar2016squad}
supplies encyclopedic paragraphs with one question each. HotpotQA~\citep{yang2018hotpotqa} is
used in its distractor configuration, contributing both supporting and distractor paragraphs;
its high document-to-question ratio is deliberate, because a corpus in which every non-poisoned
passage is topically unrelated would let a filter achieve a low false-positive rate for the
wrong reason. A small bundled toy corpus is retained so that the entire study can be reproduced
offline in seconds.

\input{\paperdir tables/tab_datasets}

\subsection{Models}

Passages and queries are embedded with \texttt{all-MiniLM-L6-v2}~\citep{reimers2019sbert,wang2020minilm},
a 384-dimensional sentence encoder, and indexed with exact inner-product
search~\citep{johnson2021faiss}. The generator is a deterministic mock model that answers
extractively from the passage most lexically aligned with the question and, unless hardened,
obeys instructions found in retrieved text. This is a deliberate instrument rather than a
convenience: it fixes the generator's behaviour exactly, so that changes in measured attack
success are attributable to the retrieval pipeline and the defense layers rather than to
sampling noise or to an opaque commercial model's undisclosed internal filtering. It models a
worst-case susceptible generator; Section~\ref{sec:res:generator} reports the same experiments
with instruction-following disabled, which bounds how much of the measured risk depends on that
assumption. Section~\ref{sec:limitations} discusses what this choice does not tell us.

\subsection{Metrics}
\label{sec:metrics}

We report attack success rate (ASR), the fraction of attacked queries meeting the adversary's
goal from Section~\ref{sec:threat}. Point estimates pool query-seed trials descriptively, and
figure error bars show the observed minimum-to-maximum range across seed-level rates. The same
questions recur across seeds, so the pooled observations are not independent. The ranges
describe variation across the experiment's random seeds; they are not confidence intervals over
a population of questions. The released result files retain only seed-level aggregates, so
query-clustered uncertainty remains unmeasured.

For suppression attacks the unconditional ASR is inflated by the pipeline's own errors: a query
the system was going to get wrong anyway is counted as a successful suppression. We therefore
also report \textbf{attributable ASR}, the success rate restricted to queries that the
\emph{same configuration} answered correctly in a clean run with no adversary present. Every
configuration is run once without an attack precisely to establish this per-configuration
baseline.

We further report clean accuracy and abstention rate (the utility cost on benign traffic);
poison retrieval rate and poison use rate (whether poison reached the candidate set, and
whether it supported the returned answer); the retrieval filter's true- and false-positive
rates at the deployed operating point, together with threshold-free AUROC and average
precision. Because poisoned passages are a small minority of the candidate set, average
precision is the more informative summary of detector quality~\citep{davis2006prc,saito2015prc}.
Ground-truth poison labels are used exclusively for scoring; no defense component reads them.
Latency is decomposed into retrieval, filtering and generation, and reported as mean and 95th
percentile.

\subsection{Configurations and Studies}

Seven configurations are compared: the undefended pipeline; each layer alone (L1--L4); the full
stack; and the full stack with signature-bound ingestion, in which a simulated trusted publisher
signs the legitimate corpus and Layer~1 requires a valid signature. Unless stated otherwise $k = 3$,
$m = 10$, $\theta_{\mathrm{trust}} = 0.6$, and the adversary's budget is one document per
target for A1--A3, five for A4 and three for A5. Every cell is repeated over three seeds, which
vary the decoy answers assigned to targets.

The studies reported in Section~\ref{sec:results} are: the main grid (configurations $\times$ attacks $\times$ corpora); the
evasion leave-one-out; the adversary budget sweep; the filter threshold sweep; the
over-retrieval sweep; the ingestion-control ladder; and the generator-hardening comparison.

\subsection{Reproducibility}

The framework, the attacks, the defense layers, the experiment suite and the code that
generates every table and figure in this paper are contained in a single repository. Results
are written incrementally as keyed JSON lines, so an interrupted sweep resumes without
recomputation, and every number reported here is regenerated from those files by a single
command rather than transcribed.

\section{Results}
\label{sec:results}

Attack-success point estimates pool query-seed trials from \nSeeds{} seeds; other summary
metrics average per-seed results. Figure error bars show the observed min-to-max range across
seeds. Unless a subsection says otherwise, figures show SQuAD and the corresponding
HotpotQA and toy-corpus panels are included where they differ materially.

\subsection{The Undefended Pipeline}

Table~\ref{tab:main} gives the full grid. Undefended, every attack succeeds. Pooled over the
four static attacks, attack success is \asrStaticSquadnone{} on SQuAD and
\asrStaticHotpotqanone{} on HotpotQA; the adaptive attack reaches \asrAdaptiveSquadnone{} and
\asrAdaptiveHotpotqanone{} respectively.

Clean accuracy of the same pipeline is \cleanAccSquadnone{} on SQuAD but only
\cleanAccHotpotqanone{} on HotpotQA. That gap is a property of the generator, not of the
defense: the mock model answers extractively from the single most query-aligned passage, and
HotpotQA questions require composing two. We report HotpotQA regardless, because its distractor
paragraphs give the retrieval filter a realistic supply of topically adjacent non-answers, but
its low benign ceiling is exactly why the attributable rate of Section~\ref{sec:metrics} is
needed: on that corpus, an unconditional suppression rate would be dominated by questions the
pipeline never answered correctly in the first place.

\input{\paperdir tables/tab_main}

\subsection{Layer Ablation Against Static Attacks}

Figure~\ref{fig:heatmap} and Figure~\ref{fig:bars} show what each layer contributes.
Ingestion verification and retrieval filtering each independently reduce pooled static attack
success on SQuAD to \casrStaticSquadLOne{} and \casrStaticSquadLTwo{} in attributable terms, and
the full stack to \casrStaticSquadfull{}. They do so by different mechanisms: Layer~1 rejects
the poison before it is indexed, so it is never retrieved, whereas Layer~2 admits it to the
index and removes it from the candidate set.

Robust aggregation alone is markedly weaker, leaving \casrStaticSquadLThree{} of static attacks
successful. Its isolate-then-aggregate design defeats prompt injection completely, because an
injected directive that reaches only one isolated generation cannot outvote the passages around
it. It does not defeat targeted poisoning, where the poisoned passage is the most query-relevant
evidence in the candidate set and therefore wins the weighted vote on its own merits. It also
carries the largest utility cost of any single layer: clean accuracy falls from
\cleanAccSquadnone{} to \cleanAccSquadLThree{} on SQuAD, because on queries where the evidence is
genuinely split the layer abstains rather than guessing.

Monitoring alone changes attack success by construction not at all (\casrStaticSquadLFour{}); it
raised alerts under every attack, which is what a detective control is for, but the ablation
makes explicit that detection is not prevention.

\begin{figure}[H]
\includegraphics[width=\textwidth]{\paperdir figures/fig2_asr_heatmap.pdf}
\caption{Attack success rate for every defense configuration and attack, on each corpus. The
rightmost column is the adaptive adversary: it is the only column in which the interior of the
grid does not change colour.\label{fig:heatmap}}
\end{figure}

\begin{figure}[H]
\includegraphics[width=\textwidth]{\paperdir figures/fig3_asr_bars.pdf}
\caption{Attack success on SQuAD; error bars show the observed min-to-max range across seeds.
Static attacks are neutralised by
retrieval filtering and by the full stack; the adaptive attack is not affected by any
content-based configuration, and only signature-bound ingestion (green) removes
it.\label{fig:bars}}
\end{figure}

\subsection{Suppression Requires the Attributable Rate}
\label{sec:res:suppression}

Context flooding, whose interaction with over-retrieval is examined in
Section~\ref{sec:res:backfill}, is the case where the unconditional rate misleads. Under the full defense its
raw ASR on SQuAD is \asrFloodingSquadfull{}, which would suggest the availability attack
half-succeeds. Restricted to queries the same configuration answered correctly with no
adversary present, its attributable ASR is \casrFloodingSquadfull{}: the flooding attack does
not, in fact, suppress a single answer that the defended pipeline would otherwise have given.
The raw figure was measuring the pipeline's own error rate, since a query the system was going
to get wrong is scored as a successful suppression. Undefended, the same attack has an
attributable ASR of \casrFloodingSquadnone{}. We therefore report both quantities throughout
and use the attributable rate whenever a suppression adversary is involved.

\subsection{The Adaptive Adversary}
\label{sec:res:adaptive}

The adaptive adversary changes the picture completely. On SQuAD its attributable success rate
is \casrAdaptiveSquadnone{} against the undefended pipeline and \casrAdaptiveSquadfull{}
against the full four-layer stack. The pooled rates show no reduction, and the observed values
agree across seeds.
Layer~2's true-positive rate against it falls to zero (Table~\ref{tab:full}): the filter does
not flag a single poisoned passage, while its false-positive rate on legitimate evidence is
unchanged. The same filter attains high detection scores against every static attack
(Table~\ref{tab:detection}), so this is not a weak detector; it is a detector whose signals the
adversary has simply declined to emit.

Layer~3 does not rescue the situation; it makes it worse. The adversary uploads several
textually distinct documents making the same claim, so the aggregation layer sees the false
answer corroborated by independent document identifiers and applies its redundancy bonus to the
adversary's group. On every corpus, enabling robust aggregation \emph{raises} adaptive attack
success relative to the undefended pipeline: from \casrAdaptiveSquadnone{} to
\casrAdaptiveSquadLThree{} on SQuAD, from \casrAdaptiveHotpotqanone{} to \casrAdaptiveHotpotqaLThree{}
on HotpotQA, and from \casrAdaptiveToynone{} to \casrAdaptiveToyLThree{} on the toy corpus, where
the effect is starkest because the true answer is supported by exactly one document and the
adversary supplies three. The mechanism that makes isolate-then-aggregate robust to a corrupted
minority is precisely the mechanism the adversary exploits: at a budget of three documents per
target, the adversary simply buys the majority. This is not a tuning failure of our particular
weighting; it is what any consistency- or corroboration-based aggregator must do when the
adversary controls more of the evidence than the corpus does.

Only the ingestion signature holds. With Layer~1 requiring a valid HMAC over the document
identifier, source identifier, and text,
the adaptive attack's attributable success falls to \casrAdaptiveSquadfullsigned{} on SQuAD and
\casrAdaptiveHotpotqafullsigned{} on HotpotQA, and clean accuracy is unchanged at
\cleanAccSquadfullsigned{}. The poisoned documents are rejected at ingestion and never reach
the index.

\input{\paperdir tables/tab_full}

\subsection{Which Evasion Technique Matters}
\label{sec:res:evasion}

Table~\ref{tab:evasion} and Figure~\ref{fig:evasion} disable one evasion technique at a time
and leave the other five enabled. Three of the six are load-bearing and three are not.

Disabling \emph{provenance spoofing} (E5) alone reduces attack success from
\evasionFullNone{} to \evasionFullSpoofsource{}: the poison is uploaded under an untrusted
source name, the allowlist rejects it, and no content filter is needed. Disabling
\emph{paraphrase} (E1) reduces it to \evasionFullParaphrase{}, because the passage then opens
with the target question and the query-echo filter recovers. Disabling the \emph{assertive
payload} (E2) reduces it to \evasionFullNodirectives{}, with the instruction-pattern filter's
true-positive rate returning to unity, since an imperative override is exactly what that
filter was built to see.

The other three techniques are, against this stack, unnecessary. Removing natural padding,
similarity dilution or cross-document corroboration leaves attack success at
\evasionFullNaturalpadding{}, \evasionFullDilutesimilarity{} and \evasionFullDiversify{}
respectively, statistically indistinguishable from the full attack. The naturalness and
similarity-outlier filters were running in all of these cells; they simply had nothing to fire
on either way. A declarative sentence about the target topic, at ordinary length, is already
fluent, already unrepetitive and already close to the corpus centroid, so padding it further
changes nothing a detector can see. E3, E4 and E6 are therefore insurance on the adversary's
side, cheap protection against detectors this deployment happens not to have, rather than
requirements against the one it does. Under signature-bound ingestion every variant fails regardless of
which content-level technique is disabled (final column of Table~\ref{tab:evasion}), because
the ingestion check does not depend on what the document says.

The middle columns of Table~\ref{tab:evasion} separate the two mechanisms by which a technique
can be load-bearing. With paraphrasing disabled the poison is still retrieved for every
attacked query, and the poison-retrieval rate stays at unity, but the filter now flags
enough of it to keep it out of the answer. With provenance spoofing disabled the poison-retrieval rate
falls to zero: nothing is filtered because nothing was ever indexed. The first is detection
after the fact, the second is prevention, and only the second leaves nothing for the adversary
to iterate against.

The practical reading is uncomfortable for content-based defense. Two of the three necessary
techniques are single rewriting decisions, namely writing a statement instead of a question
and asserting instead of commanding, that cost the adversary nothing in effectiveness. The third is a
credential problem, not a text problem.

\begin{figure}[H]
\includegraphics[width=\textwidth]{\paperdir figures/fig4_evasion.pdf}
\caption{Leave-one-out over the adaptive adversary's six evasion techniques. A bar that
collapses marks a technique the defense depends on the adversary using; under signature-bound
ingestion (green) no bar rises, because the check is independent of document
content.\label{fig:evasion}}
\end{figure}

\input{\paperdir tables/tab_evasion}

\subsection{The Ingestion-Control Ladder}
\label{sec:res:provenance}

Section~\ref{sec:res:evasion} identifies provenance spoofing as necessary for the attack of
Section~\ref{sec:res:adaptive}; Table~\ref{tab:provenance} and Figure~\ref{fig:provenance} ask which ingestion controls
actually stop it. Against a static adversary the ladder behaves as intuition suggests: adding a
trust threshold takes success from \provOpenStatic{} to \provTrustonlyStatic{}, and an
allowlist to \provAllowlistStatic{}. Against the adaptive adversary the same rungs do nothing
at all (\provOpenAdaptive{}, \provTrustonlyAdaptive{} and \provAllowlistAdaptive{}
respectively), because the source identifier and the trust score are attributes the uploader
asserts, and the
adversary asserts the ones the deployment already accepts.

The per-source quota is the one intermediate rung that does anything, and it repays a careful
reading. We set it at $1.5$ times the size of the legitimate corpus, meaning that the
contributing source may grow what it has already supplied by half. A quota expressed as a fixed
document count does not measure a defense at all: set below the corpus size it silently
truncates the knowledge base, and the resulting inability to answer is then reported as
security.

Under that policy the quota reduces adaptive attack success to \provQuotaAdaptive{} on SQuAD,
rejecting most of the adversary's uploads, and to \provQuotaAdaptiveHotpot{} on HotpotQA, where
it rejects none of them. The difference is not a property of the defense but of the ratio
between the campaign and the corpus. Attacking all $300$ SQuAD questions with three documents
each requires $900$ uploads against a remaining budget of $150$; the same campaign on HotpotQA
is $600$ uploads against a budget of nearly $1000$, and fits. A quota is a volume control, not
an authenticity control: it prices a mass campaign and is blind to a targeted one.
Section~\ref{sec:res:budget} shows the adaptive attack is already effective at one document per
target, so an adversary interested in a single high-value question never approaches any quota a
growing corpus can tolerate.

Requiring a signature takes adaptive attack success to \provSignedAdaptive{} while rejecting
every poisoned document at ingestion. This is the single clearest result in the study, and it
is not a machine-learning result: it is the observation that the distinction between a claimed
and a proved attribute, long settled in software supply-chain
security~\citep{torresarias2019intoto,newman2022sigstore,nist2022ssdf}, has not yet been
imported into retrieval pipelines.

\begin{figure}[H]
\includegraphics[width=\textwidth]{\paperdir figures/fig9_provenance.pdf}
\caption{Ingestion controls of increasing strength against a static and a provenance-spoofing
adversary. Every control that inspects a \emph{claimed} attribute is defeated. The quota rung
is a volume control rather than an authenticity control: it blunts the mass campaign measured
here and not a targeted one. Only the signature holds in both
cases.\label{fig:provenance}}
\end{figure}

\input{\paperdir tables/tab_provenance}

\subsection{Detection Quality and Operating Point}
\label{sec:res:detection}

Table~\ref{tab:detection} reports threshold-free detector quality and Figure~\ref{fig:roc} the
empirical ROC traced by sweeping the calibrated suspicion threshold $\tau$ of
Equation~(\ref{eq:suspicion}). Both are measured in the filter-only configuration, because in
the full stack Layer~1 rejects the static poison before it is ever a retrieval candidate and
the filter's score would otherwise be credited with Layer~1's work.

Against the static attacks the filter is close to separable: AUROC is at least
\aurocStaticMinSquad{} on SQuAD, \aurocStaticMinHotpotqa{} on HotpotQA and
\aurocStaticMinToy{} on the toy corpus, with average precision to match. Against the adaptive
attack AUROC is \aurocAdaptiveSquad{}, \aurocAdaptiveHotpotqa{} and \aurocAdaptiveToy{}
respectively. These sit at chance, and on two of the three corpora below it, meaning the
suspicion score
ranks adaptive poison as slightly \emph{less} suspicious than the legitimate passages it
competes with. At the deployed operating point the corresponding true-positive rates are
\tprAdaptiveSquad{}, \tprAdaptiveHotpotqa{} and \tprAdaptiveToy{}.

This is worth stating plainly, because a detector reported only at its default threshold and
only against published attack constructions would look excellent.

Figure~\ref{fig:threshold} shows what moving $\tau$ costs. Detection of the adaptive attack
\emph{can} be bought, but not at a price any deployment would pay. At the default operating
point $\tau = \tauDefault{}$ the filter catches \tauDefaultTprAdaptive{} of adaptive poison at a
false-positive rate of \tauDefaultFpr{}, leaving attack success at \tauDefaultAsrAdaptive{} with
clean accuracy \tauDefaultCleanAcc{}. Tightening to $\tau = \tauModerate{}$ raises adaptive
detection to \tauModerateTprAdaptive{} and cuts attack success to \tauModerateAsrAdaptive{}, but
clean accuracy has already collapsed to \tauModerateCleanAcc{}, because the filter is now
discarding
\tauModerateFpr{} of legitimate evidence. At $\tau = \tauAggressive{}$ attack success is
\tauAggressiveAsrAdaptive{}, and clean accuracy is \tauAggressiveCleanAcc{}: the system has been
secured by being switched off. Throughout the sweep the static attacks are caught at a
true-positive rate of essentially one, so no single threshold serves both adversaries.

The reason is that the adaptive passages are not merely sitting below the threshold. They are
indistinguishable on these features from genuine encyclopedic prose, which is what they were
written to be, so the only threshold that captures them is one that captures the corpus too.

\input{\paperdir tables/tab_detection}

\begin{figure}[H]
\includegraphics[width=\textwidth]{\paperdir figures/fig6_detection_roc.pdf}
\caption{Empirical ROC of the retrieval filter, traced by sweeping the calibrated suspicion
threshold. The adaptive attack (black) lies on the chance diagonal on every
corpus.\label{fig:roc}}
\end{figure}

\begin{figure}[H]
\includegraphics[width=0.62\textwidth]{\paperdir figures/fig7_threshold.pdf}
\caption{The retrieval filter's operating point on SQuAD. The static constructions are detected
at a true-positive rate of essentially one across the whole sweep; the adaptive attack is
detected only once the threshold is low enough to discard a large share of legitimate evidence
as well, at which point clean accuracy has already collapsed. No single threshold serves both
adversaries.\label{fig:threshold}}
\end{figure}

\begin{figure}[H]
\includegraphics[width=\textwidth]{\paperdir figures/fig8_backfill.pdf}
\caption{Over-retrieval at two operating points. At the default threshold (left) the filter
discards almost nothing, so raising the candidate multiplier changes no outcome and only buys
latency. Under aggressive filtering (right) backfill repairs availability, since abstention
and attributable attack success both fall sharply with $m$, but recovers only a fraction of the
lost clean accuracy, because evidence discarded by a mistuned filter cannot be
replaced.\label{fig:backfill}}
\end{figure}

\subsection{Adversary Budget}
\label{sec:res:budget}

Figure~\ref{fig:budget} sweeps the number of documents $b$ the adversary inserts per target.
The undefended pipeline is saturated at a single document, consistent with the retrieval
condition PoisonedRAG identifies: one passage constructed to embed next to the query is enough
to enter a top-3 context. Under the full defense the static curve stays near zero across the
whole range, so more poison does not help an adversary whose construction is detectable. Under
signature-bound ingestion the adaptive curve is flat at zero: budget is irrelevant when every
document is rejected before indexing.

The adaptive curve against the full defense has a shape worth reading closely, because it
prices the attack. Against the \emph{undefended} pipeline it is flat, at
\budgetAdaptiveNoneSquadBOne{} for $b=1$ and \budgetAdaptiveNoneSquadBTwenty{} for $b=20$ on SQuAD,
since one retrievable passage already carries the false claim. Against the
\emph{defended} pipeline it steps: \budgetAdaptiveFullSquadBOne{} at $b=1$, rising to
\budgetAdaptiveFullSquadBTwo{} at $b=2$ and flat thereafter. The step is starkest on the toy
corpus, from \budgetAdaptiveFullToyBOne{} to \budgetAdaptiveFullToyBTwo{}.

That step is Layer~3's redundancy bonus changing hands. With one poisoned document the
adversary's claim is a minority of one against the corpus's genuine passage and the aggregator
still resolves in favour of the truth a meaningful share of the time; with two, the false claim
carries the corroboration bonus and the genuine passage does not. The adversary's price for
turning the aggregation layer from an obstacle into an instrument is therefore \emph{one
additional document}. Everything above $b=2$ is wasted effort, which is precisely the wrong
shape for a defender: the attack that works leaves the smallest footprint, and volume-based
monitoring has nothing to see.

\begin{figure}[H]
\includegraphics[width=\textwidth]{\paperdir figures/fig5_budget.pdf}
\caption{Attack success against the adversary's per-target upload budget. The adaptive attack
is effective from the smallest budget, so its footprint offers little for volume-based
monitoring to detect.\label{fig:budget}}
\end{figure}

\subsection{Over-Retrieval, Utility and Latency}
\label{sec:res:backfill}

It is often argued, and we argued it ourselves when designing
Algorithm~\ref{alg:pipeline}, that a retrieval filter must be paired with over-retrieval, since dropping a candidate otherwise just
shortens the context and hands a suppression adversary a win for free. Figure~\ref{fig:backfill}
tests that argument by sweeping the candidate multiplier $m$ at two operating points, and the
answer depends entirely on which one is deployed.

At the default threshold (left panel) the argument is simply void. Going from $m=1$ to $m=20$
leaves clean accuracy at \backfillDefaultAccLo{}, abstention at \backfillDefaultAbstainLo{} and
attributable static ASR at \backfillDefaultAsrLo{}, identical to two decimal places, while
p95
latency rises from \backfillDefaultLatLo{}~ms to \backfillDefaultLatHi{}~ms. Nothing is
backfilled because nothing is dropped: with ingestion control in place the filter discards under
one percent of candidates (\backfillDefaultFprHi{}), so over-retrieval is pure cost. A
deployment that sets a large multiplier for safety here is paying an order of magnitude in
latency for no security and no utility.

Under aggressive filtering (right panel, $\tau = 0.6$, false-positive rate
\backfillStrictFprLo{} at $m=1$) the argument holds, but only for part of what it claims.
Backfill repairs the \emph{availability} damage convincingly: abstention falls from
\backfillStrictAbstainLo{} to \backfillStrictAbstainHi{} and attributable static ASR from
\backfillStrictAsrLo{} to \backfillStrictAsrHi{} as $m$ grows. It does not repair
\emph{utility}: clean accuracy recovers only from \backfillStrictAccLo{} to
\backfillStrictAccHi{}, still far below the \backfillDefaultAccHi{} obtained at the default
threshold, because even at $m=20$ the filter is discarding \backfillStrictFprHi{} of legitimate
evidence and no amount of over-retrieval manufactures evidence that was thrown away.

The practical conclusion is that backfill is a necessary companion to aggressive filtering and
not a substitute for choosing the threshold correctly. The multiplier is the wrong lever: it can
stop a badly-tuned filter from denying service, but it cannot make a badly-tuned filter useful.

Table~\ref{tab:utility} reports the utility cost of each configuration on unattacked traffic and
Table~\ref{tab:latency} the latency decomposition. The largest utility cost is Layer~3, which
takes clean accuracy from \cleanAccSquadnone{} to \cleanAccSquadLThree{} on SQuAD and raises the
abstention rate from near zero to roughly a third of queries: on questions where the isolated
answers disagree it declines to answer rather than picking one.

The measured time is dominated by filtering, at roughly \filterMsFullSquad{}~ms of a
\meanMsFullSquad{}~ms defended query, because the layer scores $mk = 30$ candidates with four
detectors, one of which maintains corpus-level statistics. This ranking is an artefact of the
harness and should not be carried into a deployment: the generator here is a local mock costing
a fraction of a millisecond per call, so the one component that would dominate a real system is
precisely the one this measurement makes invisible. The structural cost of Layer~3 is not a
matter of measurement at all: by construction it issues one generator call per retained
passage instead of one per query, so at $k=3$ it triples generator spend regardless of how fast
the generator is. The honest reading of Table~\ref{tab:latency} is therefore that the defense's
\emph{own} overhead is dominated by candidate filtering and scales with $m$, while its dominant
\emph{deployment} cost is the $k$-fold multiplication of generator calls.

\input{\paperdir tables/tab_utility}
\input{\paperdir tables/tab_latency}

\begin{figure}[H]
\includegraphics[width=0.62\textwidth]{\paperdir figures/fig10_latency.pdf}
\caption{Where the defended pipeline spends its time. Filtering is inexpensive; the cost is
Layer~3, which issues one generator call per retained passage instead of one per
query.\label{fig:latency}}
\end{figure}

\begin{figure}[H]
\includegraphics[width=\textwidth]{\paperdir figures/fig11_frontier.pdf}
\caption{Security--utility frontier. The horizontal axis is clean accuracy, the vertical axis
one minus attack-attributable success, so the top-right corner is best. Against the static
attacks (circles) the defended configurations sit at the top of every panel. Against the
adaptive attack (crosses) the same configurations fall to the bottom, with the exception of
signature-bound
ingestion, whose cross stays at the top at no additional utility cost. Points are offset
horizontally by configuration so that coincident values remain
visible.\label{fig:frontier}}
\end{figure}

\subsection{How Much Depends on a Susceptible Generator}
\label{sec:res:generator}

The generator used throughout obeys instructions found in retrieved text, which is a worst-case
assumption. Table~\ref{tab:generator} and Figure~\ref{fig:generator} repeat the experiments
with that behaviour disabled. Prompt injection is, as expected, eliminated: with no
instruction-following there is no instruction hierarchy to override. The other four attacks are
essentially unaffected, including the adaptive attack, because none of them relies on the
generator obeying an order. They rely on the generator believing a document, which any
faithful, grounded generator will do by design. Indeed, a generator that ignored its
retrieved context would defeat the purpose of retrieval. Instruction-hierarchy training and structured
prompting~\citep{wallace2024hierarchy,chen2024struq,chen2025secalign} therefore address one of
the five attacks in this study and leave the knowledge-integrity problem untouched.

\input{\paperdir tables/tab_generator}

\begin{figure}[H]
\includegraphics[width=0.62\textwidth]{\paperdir figures/fig12_generator.pdf}
\caption{Hardening the generator against injected instructions removes prompt injection and
leaves the poisoning attacks essentially unchanged: they exploit grounding, not
obedience.\label{fig:generator}}
\end{figure}

\section{Discussion}
\label{sec:discussion}

\subsection{What the Layers Are Actually Doing}

Read together, the ablation of Section~\ref{sec:results} and the evasion study give each
layer a more precise role than ``defense in depth'' usually implies, and one that the
suppression correction of Section~\ref{sec:res:suppression} is needed to see clearly.

Layer~1 is the only \emph{preventive} control in the stack that an adaptive adversary cannot
argue with, but only in its cryptographic form. Its allowlist and trust threshold check
assertions, and an adversary who can write to the corpus can assert anything. This distinction
is invisible against static attacks, which conveniently upload from an untrusted source name,
and it is the whole result against an adaptive one.

Layer~2 is a \emph{detective} control that has been widely deployed as though it were
preventive. It is genuinely effective against the published attack constructions, and it
generalises to nothing beyond them: its four signals correspond to four things an adversary can
simply not do. We do not conclude that retrieval filtering is useless, since it raises the
cost of the naive attack and produces the evidence an incident responder needs, but it should
be
budgeted as detection, and a system whose security case rests on it does not have one.

Layer~3 (Section~\ref{sec:l3}) is the layer whose behaviour most needs restating. Against prompt injection it is the
strongest single control in the stack, and its abstention is a genuinely useful failure mode:
under attack it prefers to answer nothing over answering wrongly. But its robustness rests on
the corrupted passages being a minority, and an adversary with a budget of three documents per
target owns that majority. Our measurements show this is not merely a missed opportunity: on
every corpus, enabling robust aggregation \emph{increased} adaptive attack success, because the
redundancy bonus that rewards corroboration rewards whoever supplies the most of it. It is also
the most expensive layer in both utility (Figure~\ref{fig:frontier}) and compute
(Figure~\ref{fig:latency}). Certified variants make the underlying assumption
explicit~\citep{xiang2024robustrag}, which is the right response: the guarantee is conditional
on a corruption bound, and a deployment has to enforce that bound somewhere else, namely at
ingestion.
A corroboration-weighted aggregator deployed without such a bound is not a neutral addition; it
is an amplifier pointed in whichever direction the evidence majority lies.

Layer~4 does what detective controls do. It alerted under every attack, and it cannot stop any
of them. Its value is operational, and our ablation is included precisely so that this is not
confused with prevention.

\subsection{Guidance for Deployments}

Four recommendations follow from the measurements rather than from the architecture.

\emph{Sign at the source, verify at ingestion} (Section~\ref{sec:res:provenance}). Bind each document to a key held by the
ingestion pipeline, not to a source name in a metadata field. This is the only control in our
study that survived, it costs one HMAC per document, and its clean accuracy cost was nil.
Where content genuinely must be accepted from unauthenticated sources, treat that corpus as a
separate, lower-trust index rather than mixing it with signed content.

\emph{Budget content filters as detection.} Deploy them, alert on them, retain their scores for
investigation, and do not let them appear in the security case as a preventive control. Report
their performance against rewritten payloads, not only against published constructions.

\emph{Tune the filter's threshold; do not compensate with over-retrieval.} The two knobs are
not interchangeable, and we found the intuitive advice to be wrong in the regime deployments
actually run in. At a sensible threshold the filter discards so little that raising the
candidate multiplier changes no outcome at all and multiplies latency roughly tenfold. At an
aggressive threshold, over-retrieval does prevent the filter from denying service, since
abstention and attack success both fall steeply with $m$, but it recovers only a fraction of
the clean
accuracy the threshold gave away. Set the operating point from a measured false-positive rate
on legitimate traffic, use the smallest multiplier that keeps abstention flat, and treat a
large multiplier as a symptom of a mistuned filter rather than a safety margin.

\emph{Do not conflate instruction-hierarchy defenses with knowledge integrity.} Hardening the
generator against injected instructions eliminated exactly one of our five attacks. The rest
attack what the model is given to believe, and a well-behaved grounded generator will believe
it.

\subsection{Why Content-Based Detection Fails Here}

The failure is structural rather than a matter of tuning. Each of our detectors keys on an
artefact of a \emph{convenient} attack construction: the query is echoed because that is the
easiest way to guarantee retrieval; the payload is imperative because that is the easiest way
to steer a generator; the text is repetitive because it was templated. None of these is
required. A passage that asserts a false fact in fluent prose, on topic, at ordinary length, is
by construction indistinguishable from a passage that asserts a true fact in fluent prose.
The difference is the truth of the claim, which is not a property of the text's surface form. This
is why the adaptive ROC in Figure~\ref{fig:roc} sits on the diagonal, and why we expect
detectors built on richer surface features, including perplexity-based
ones~\citep{alon2023perplexity,jain2023baseline}, to face the same ceiling. Distinguishing true
from false claims is a knowledge problem, and the pipeline's answer to knowledge problems is to
consult the corpus, which the adversary controls.

\subsection{Relation to Prior Results}

Our static-attack results reproduce the qualitative findings of the poisoning literature: a
single crafted passage is enough to hijack a targeted query in an undefended
pipeline~\citep{zou2024,zhong2023}, and availability-style poisoning suppresses answers without
substituting one~\citep{shafran2024jamming}. Our defended results are consistent with reported
gains for retrieval-time detection and isolate-then-aggregate
generation~\citep{xiang2024robustrag,zhou2025trustrag} \emph{when evaluated the same way those
works evaluate}. The divergence appears only under the adaptive adversary, which is the
predicted consequence of the methodological gap that
\citet{carlini2017bypassing,athalye2018obfuscated,tramer2020adaptive} documented in a different
subfield. We read our results not as a refutation of those defenses but as evidence that the
RAG security literature is at the point in its development where adaptive evaluation must
become standard.

\section{Limitations}
\label{sec:limitations}

The generator is a deterministic mock. This is deliberate, since it isolates the retrieval
pipeline and makes attack success measurable without an opaque model in the loop, but it is not
a production language model. A stronger model might resist a single false assertion that
contradicts its parametric knowledge, and the generator study bounds only the
instruction-following component of that question. Replicating this study across current hosted
and open-weight models is the most important next step, and we expect the adaptive attack's
advantage to shrink but not vanish, since a model asked to ground its answer in a corroborated
retrieved claim has little basis for refusing.

Our corpora range from 30 to 1,999 documents, far below the millions in an enterprise index. Scale affects
the similarity-outlier filter, which estimates corpus statistics, and it affects how easily a
poisoned passage reaches the top-$k$ for a query. It does not obviously affect the central
result, since provenance spoofing is scale-independent, but this is an assumption rather than a
measurement.

The adaptive adversary is one adversary, constructed by us against our own defense. It gives a
lower bound on what an adaptive attacker achieves, not an upper bound, and a stronger one
almost certainly exists, in particular one that optimises passage embeddings against the
retriever directly rather than relying on lexical overlap. Symmetrically, our defense is one
implementation; a deployment with different filters would need its own adaptive evaluation,
which is the point.

Attacks and defenses were implemented in the same codebase, which risks benchmark coupling. We
mitigated this by exposing evasion techniques as independent switches and reporting per-layer
attribution rather than aggregate scores, so a reader can see which defensive assumption each
result depends on, but independent reimplementation would be stronger evidence.

The signature result assumes the signing key and signing authority are protected and only
authorized publishers can request signatures. The harness uses a fresh in-memory HMAC key and
does not test compromise of the signer or poisoning by a malicious authorized publisher. The
signature binds a versioned canonical envelope containing document text, document identifier,
and source identifier, but does not enforce freshness or prevent replay of an old signed
document. A deployment still needs signer authorization, replay controls, key rotation, and
revocation.

Finally, reported latencies exclude external model and network cost, and the abstention
behaviour we report is a property of our aggregation thresholds rather than a calibrated
operating point that a deployment should adopt unexamined.

\section{Conclusions}

We built a four-layer defense against knowledge-base poisoning in retrieval-augmented
generation (ingestion verification, retrieval filtering, robust aggregation and anomaly
monitoring) and evaluated it on three corpora against five attacks, including a defense-aware
adaptive adversary of our own construction.

Against the four static attacks from the literature the architecture performs as expected:
attack success falls from near unity to approximately zero in attributable terms, at a modest
utility cost. Against an adversary who merely rewrites the payload, making it declarative
instead of interrogative and assertive instead of imperative, and uploading it under a source
name the deployment already trusts, the same architecture fails. The retrieval filter's true-positive rate falls
to zero and its ROC to the chance diagonal; robust aggregation is turned against itself by
manufactured corroboration; and attack success is statistically indistinguishable from the
undefended pipeline. A leave-one-out ablation shows the adversary needs exactly three of its six
techniques, two of them free rewriting choices.

Requiring an unforgeable signature at ingestion restored the defense completely, at no cost to
clean accuracy, and was the only control in the study that did. We take this as the practical
conclusion: content-based heuristics belong in a RAG security architecture as detection and as
incident evidence, but the preventive control has to be verifiable provenance. The security
question for a retrieval pipeline is not primarily what a document says, since an adversary
chooses that, but whether the pipeline can prove who put it there.
